\chapter{Coppie ordinate e corrispondenze}\label{cap:corrispondenze}

Gli insiemi non hanno un ordine: $\{x, y\} = \{y, x\}$. Per parlare di funzioni e relazioni ci servono invece
coppie in cui conta \textbf{chi viene prima}. In questo capitolo costruiamo le coppie ordinate usando solo gli
assiomi, poi il prodotto cartesiano e infine le corrispondenze tra insiemi.

% ──────────────────────────────────────────────────────────────
\section{Coppie ordinate}\label{sec:coppie}

\dfn[coppia-ordinata]{Coppia ordinata}{
    Siano $x$ e $y$ insiemi. L'insieme $\{\{x\}, \{x, y\}\}$ si chiama \textbf{coppia ordinata} di
    \textbf{prima coordinata} $x$ e \textbf{seconda coordinata} $y$, e si scrive $(x, y)$.
}

L'insieme esiste per l'assioma della Coppia, applicato prima a $x, x$ e $x, y$ e poi a $\{x\}, \{x, y\}$.
La proprietà che ci interessa è che due coppie ordinate sono uguali solo se hanno le stesse coordinate
\textbf{nello stesso ordine}.

\begin{concetto}
\begin{Teorema}{Proprietà caratteristica della coppia ordinata}{coppia-ordinata}
    \begin{prerequisiti}
        \dip{Definizione}{def:coppia-ordinata}
        \dip{Assioma}{ass:estensionalita}
        \dip{Assioma}{ass:coppia}
        \dip{Definizione}{def:notazioni-assiomi}
        \dip{Proposizione}{prop:doppiaimp}
    \end{prerequisiti}
    \[ \forall x, y, w, z\,\big((x, y) = (w, z) \leftrightarrow (x = w \wedge y = z)\big) \]
\end{Teorema}
\begin{dimostrazione}
    $(\leftarrow)$ Se $x = w$ e $y = z$, allora ovviamente $(x, y) = (w, z)$.

    $(\rightarrow)$ Supponiamo $(x, y) = (w, z)$, cioè $\{\{x\}, \{x, y\}\} = \{\{w\}, \{w, z\}\}$.
    \begin{itemize}
        \item \textbf{Se $x = y$}, allora $\{\{x\}, \{x, x\}\} = \{\{x\}\}$ ha un solo elemento, quindi anche
              $\{\{w\}, \{w, z\}\}$ ne ha uno solo: $\{w\} = \{w, z\} = \{x\}$. Quindi $w = x$ e $z = w$, cioè $x = y = w = z$.
        \item \textbf{Se $x \neq y$}, allora $\{x\} \neq \{x, y\}$ e $\{x, y\}$ ha due elementi. Per Estensionalità
              $\{x, y\}$ deve essere uguale a $\{w\}$ oppure a $\{w, z\}$; non può essere $\{w\}$ (altrimenti $x = w = y$),
              quindi $\{x, y\} = \{w, z\}$. Allora l'altro elemento, $\{x\}$, deve essere uguale a $\{w\}$, cioè $x = w$.
              Ma allora $\{x, y\} = \{x, z\}$ e $y$, non potendo essere uguale a $x$, deve essere uguale a $z$.
    \end{itemize}
    In entrambi i casi $x = w$ e $y = z$.
\end{dimostrazione}
\end{concetto}

\dfn[terna-ordinata]{Terna ordinata}{
    Siano $x$, $y$, $z$ insiemi. La \textbf{terna ordinata} di coordinate $x$, $y$, $z$ è
    \[ (x, y, z) := \big((x, y), z\big). \]
}

\begin{concetto}
\begin{Corollario}{Proprietà caratteristica della terna ordinata}{terna}
    \begin{prerequisiti}
        \dip{Teorema}{teo:coppia-ordinata}
        \dip{Definizione}{def:terna-ordinata}
    \end{prerequisiti}
    $(x, y, z) = (x', y', z') \leftrightarrow (x = x' \wedge y = y' \wedge z = z')$.
\end{Corollario}
\begin{dimostrazione}
    Per il Teorema~\ref{teo:coppia-ordinata} applicato due volte:
    $\big((x,y),z\big) = \big((x',y'),z'\big) \leftrightarrow (x,y) = (x',y') \wedge z = z' \leftrightarrow x = x' \wedge y = y' \wedge z = z'$.
\end{dimostrazione}
\end{concetto}

% ──────────────────────────────────────────────────────────────
\section{Prodotto cartesiano}\label{sec:prodotto-cartesiano}

Vogliamo l'insieme di \textbf{tutte} le coppie ordinate con prima coordinata in $a$ e seconda coordinata in $b$.
Per costruirlo con la Separazione ci serve un insieme che le contenga già tutte.

\begin{concetto}
\begin{Lemma}{Dove vivono le coppie ordinate}{coppia-parti}
    \begin{prerequisiti}
        \dip{Definizione}{def:coppia-ordinata}
        \dip{Assioma}{ass:parti}
        \dip{Definizione}{def:inclusione}
        \dip{Proposizione}{prop:appartenenza}
    \end{prerequisiti}
    Siano $a$, $b$ insiemi, $x \in a$ e $y \in b$. Allora $(x, y) \in \mcP(\mcP(a \cup b))$.
\end{Lemma}
\begin{dimostrazione}
    \textbf{Primo insieme delle parti.} Dato che $x \in a \cup b$, il singleton $\{x\}$ è un sottoinsieme di $a \cup b$,
    quindi $\{x\} \in \mcP(a \cup b)$. Allo stesso modo, poiché $x, y \in a \cup b$, anche $\{x, y\} \in \mcP(a \cup b)$.

    \textbf{Secondo insieme delle parti.} Poiché sia $\{x\}$ sia $\{x, y\}$ appartengono a $\mcP(a \cup b)$, l'insieme che
    li racchiude entrambi, $\{\{x\}, \{x, y\}\} = (x, y)$, è un sottoinsieme di $\mcP(a \cup b)$.
    Quindi $(x, y) \in \mcP(\mcP(a \cup b))$.
\end{dimostrazione}
\end{concetto}

\dfn[prodotto-cartesiano]{Prodotto cartesiano}{
    Siano $a$ e $b$ insiemi. Il \textbf{prodotto cartesiano} di $a$ e $b$ è
    \[ a \times b := \{z \in \mcP(\mcP(a \cup b)) \mid \exists x, y\,(x \in a \wedge y \in b \wedge z = (x, y))\}. \]
    Per il Lemma~\ref{lem:coppia-parti}, $a \times b$ contiene \textbf{tutte} le coppie $(x, y)$ con $x \in a$ e $y \in b$.
}

\info{Perché le parti delle parti?}{
    Facciamo le parti delle parti di $a \cup b$ perché, partendo dalla definizione di coppia
    $\{\{x\}, \{x, y\}\}$, sappiamo solo che $x \in a$ e $y \in b$: di conseguenza sia $x$ sia $y$ appartengono
    all'insieme unione $a \cup b$. È il più piccolo insieme «già costruito» in cui cercare le coppie con la Separazione.
}

% ──────────────────────────────────────────────────────────────
\section{Corrispondenze}\label{sec:corrispondenze}

\dfn[corrispondenza]{Corrispondenza}{
    Siano $a$, $b$, $g$ insiemi con $g \subseteq a \times b$. Allora la terna $\rho = (a, b, g)$ si dice
    \textbf{corrispondenza} tra $a$ e $b$, e $g$ si dice \textbf{grafico} della corrispondenza.
    L'insieme delle corrispondenze tra $a$ e $b$ lo chiamiamo $\mathrm{Corr}(a, b)$.

    Se $\rho = (a, b, g) \in \mathrm{Corr}(a, b)$, $x \in a$, $y \in b$ e $(x, y) \in g$, scriveremo $x \rho y$.
}

\dfn[relazione-binaria]{Relazione binaria}{
    Una corrispondenza $(a, a, g)$ di un insieme con se stesso si dice \textbf{relazione binaria} su $a$.
    Definiamo $\mathrm{Rel}(a) := \mathrm{Corr}(a, a)$ e la \textbf{diagonale} $\mathrm{diag}(a) := \{(x, x) \mid x \in a\}$.
    \begin{itemize}
        \item Se $g = a \times a$, la relazione si dice \textbf{totale}.
        \item Se $g = \mathrm{diag}(a)$, la relazione si dice \textbf{identica}.
    \end{itemize}
}

\dfn[prodotto-corrispondenze]{Prodotto tra corrispondenze}{
    Siano $\rho = (a, b, g_1)$ e $\sigma = (c, d, g_2)$ corrispondenze. Definisco la corrispondenza
    \[ \sigma\rho = (a, d, g_3) \]
    e la chiamo \textbf{prodotto tra $\rho$ e $\sigma$}, dove
    \[ \forall x, y\,\Big((x, y) \in g_3 \leftrightarrow \exists z\,\big((x, z) \in g_1 \wedge (z, y) \in g_2\big)\Big). \]
}

Attenzione all'ordine: in $\sigma\rho$ si applica \textbf{prima} $\rho$ e \textbf{poi} $\sigma$. Le coppie di $g_3$ sono quelle
«collegate attraverso elementi in comune tra $b$ e $c$».

\begin{esempio}{Prodotto tra corrispondenze}
    Sia $\rho$ da $a = \{x_1, x_2\}$ a $b = \{z_1, z_2\}$ e $\sigma$ da $c = b$ a $d = \{y_1, y_2\}$:
    \begin{center}
    \begin{tikzpicture}[>=Stealth, every node/.style={font=\small}]
        \node (x1) at (0,0.6) {$x_1$}; \node (x2) at (0,-0.6) {$x_2$};
        \node (z1) at (3,0.6) {$z_1$}; \node (z2) at (3,-0.6) {$z_2$};
        \node (y1) at (6,0.6) {$y_1$}; \node (y2) at (6,-0.6) {$y_2$};
        \draw[rounded corners] (-0.5,-1.1) rectangle (0.5,1.1); \node at (0,1.4) {$a$};
        \draw[rounded corners] (2.5,-1.1) rectangle (3.5,1.1); \node at (3,1.4) {$b = c$};
        \draw[rounded corners] (5.5,-1.1) rectangle (6.5,1.1); \node at (6,1.4) {$d$};
        \draw[->,primary] (x1) -- (z1); \draw[->,primary] (x1) -- (z2);
        \draw[->,cherryred] (z1) -- (y2);
        \node[primary] at (1.5,-1.4) {$\rho$}; \node[cherryred] at (4.5,-1.4) {$\sigma$};
    \end{tikzpicture}
    \end{center}
    Qui $x_1 \rho z_1$ e $z_1 \sigma y_2$, quindi $x_1 \,\sigma\rho\, y_2$. Nessun'altra coppia è collegata:
    il grafico di $\sigma\rho$ è $\{(x_1, y_2)\}$.
\end{esempio}

\begin{esempio}{Altri prodotti}
    \begin{itemize}
        \item Siano $\rho$ e $\sigma$ le relazioni binarie tra gli esseri umani «essere padre di» ed «essere fratello di».
              Allora $x \,\rho\sigma\, y$ se e solo se «$x$ è padre del fratello di $y$», $x \,\sigma\rho\, y$ se e solo se
              «$x$ è fratello del padre (cioè è uno zio) di $y$», mentre $x \rho^2 y$ (con $\rho^2 = \rho\rho$) se e solo se
              «$x$ è il nonno paterno di $y$».
        \item Consideriamo l'inclusione su $\mcP(a)$. Allora $\subseteq^2 = \subseteq$, grazie alla transitività dell'inclusione
              (Proposizione~\ref{prop:transitivita-inclusione}).
    \end{itemize}
\end{esempio}

\begin{concetto}
\begin{Proposizione}{Associatività del prodotto tra corrispondenze}{assoc-prodotto}
    \begin{prerequisiti}
        \dip{Definizione}{def:prodotto-corrispondenze}
        \dip{Corollario}{cor:terna}
        \dip{Assioma}{ass:estensionalita}
        \dip{Proposizione}{prop:associativita}
    \end{prerequisiti}
    Il prodotto di corrispondenze è \textbf{associativo}: per ogni terna di corrispondenze $\rho$, $\sigma$, $\tau$
    \[ \tau(\sigma\rho) = (\tau\sigma)\rho. \]
\end{Proposizione}
\begin{dimostrazione}
    Siano $\rho = (a, b, g_1)$, $\sigma = (c, d, g_2)$, $\tau = (e, h, g_3)$. Entrambi i prodotti hanno $a$ come primo
    insieme e $h$ come secondo, quindi (Corollario~\ref{cor:terna}) basta mostrare che i grafici coincidono.
    Per definizione di prodotto, applicata due volte:
    \[
    \begin{aligned}
        (x, y) \in \text{grafico di } \tau(\sigma\rho) &\leftrightarrow \exists w\,\Big(\exists z\,\big((x, z) \in g_1 \wedge (z, w) \in g_2\big) \wedge (w, y) \in g_3\Big)\\
        &\leftrightarrow \exists z\,\exists w\,\big((x, z) \in g_1 \wedge (z, w) \in g_2 \wedge (w, y) \in g_3\big),\\
        (x, y) \in \text{grafico di } (\tau\sigma)\rho &\leftrightarrow \exists z\,\Big((x, z) \in g_1 \wedge \exists w\,\big((z, w) \in g_2 \wedge (w, y) \in g_3\big)\Big)\\
        &\leftrightarrow \exists z\,\exists w\,\big((x, z) \in g_1 \wedge (z, w) \in g_2 \wedge (w, y) \in g_3\big).
    \end{aligned}
    \]
    Nei passaggi abbiamo portato fuori il quantificatore $\exists$ (la variabile quantificata non compare nell'altro
    membro della congiunzione) e usato l'associatività di $\wedge$. Le due condizioni coincidono, quindi per
    Estensionalità i grafici sono uguali.
\end{dimostrazione}
\end{concetto}
